\section{Introduction}
Embodied AI for aerospace vehicles links onboard sensing, task reasoning, and executable actions under vehicle and mission constraints. In aerial disaster inspection, an agent may be asked whether a building is damaged, how many damaged buildings lie within a region, or where the nearest damaged building is located. Descending toward a building may improve local detail while narrowing the visible footprint and excluding evidence needed for a regional assessment. The value of an additional observation therefore depends on the question being answered and the maneuver it requires.

Embodied multimodal models connect visual and state observations to task reasoning, while language-guided navigation and vision-language-action systems ground plans in executable skills and actions~\cite{driess2023palme,shah2023lmnav,ichter2023saycan,zitkovich2023rt2}. Aerial navigation and embodied benchmarks evaluate action-conditioned observations and mission behavior; remote-sensing agents support disaster interpretation and structured tool use~\cite{liu2023aerialvln,lee2025citynav,gao2026openfly,yao2024aeroverse,guo2026bedi,zhao2025cityeqa,liu2024rescueadi,thinkgeo2025}. These capabilities motivate an evaluation question that is narrower but consequential: how does a vehicle's decision to acquire a new view affect the answer to a task whose geographic scope and reference remain fixed? Answering it requires a controlled observation change and a trace linking the action to geographic evidence and the answer call. A standardized episode interface can make this comparison repeatable by distinguishing a requested action from an executed movement and separating changed-view effects from another answer opportunity.

We present DisasterClaw, an imagery-based platform and task benchmark for this evaluation problem. Paired xBD tiles, building annotations, task regions, and simulated vehicle states share a common geographic reference. Annotation-derived answers remain fixed as the agent searches or changes its observation footprint. Four task families impose different requirements for local detail, regional coverage, and geographic reference. An operator console supports scene inspection, and a headless runner executes the quantitative evaluation.

The evaluation connects observation changes to building predictions and task answers. A camera-ladder diagnostic on the development set measures how a narrower footprint affects building localization and damage classification, including corrections and harms. On the frozen final set of 160 human-reviewed questions from 44 regions in three events, shortcut controls assess performance using only question text or permitted runtime metadata. A fixed-view matched-call diagnostic then branches at every executed T1 recheck: one branch reruns perception and answering on the previous image, and the other executes the logged recheck. Matching the answer opportunity helps distinguish the effect of changed observations at individual transitions. A separate MESSI case compares real UAV images acquired at two heights with repeat-view and digital-crop controls. We then inventory the full 2,525-image release and select 409 eligible two-height vegetation regions under the fixed protocol. On these pairs, we run a single-temporal adaptation of DisasterClaw with its existing YOLO and SegFormer models; the controller sees only the high view until its recheck gate authorizes the registered lower view. This measures adapted offline re-observation on real UAV imagery, alongside the platform's primary xBD evaluation.

The contributions are:
\begin{enumerate}
\item a georeferenced, ROI-scoped task benchmark in which the reference answer stays fixed as the observation changes, with an explicit boundary between runtime task cues and scoring labels;
\item an inspection interface that records requested and executed actions, rendered views, predicted evidence, answer calls, and terminal outcomes within the same episode; and
\item an empirical characterization of observation value through shortcut controls, fixed-view matched calls, and corrections and harms by task, using a fixed task-conditioned policy as a reference implementation.
\end{enumerate}

The platform renders existing orthographic imagery under simplified motion. We position it as an evaluation component for embodied aerial agents: it measures task-conditioned sensing decisions and answer consequences, while leaving vehicle dynamics, onboard sensing, and physical flight to future evaluation. The findings apply to this controlled observation setting; they do not establish flight effectiveness.
